\section{Limitations and Future Work}
\label{sec:limitations}
 
We discuss the current limitations of Omni-Shield
honestly, alongside concrete directions for future work.
 
% ─────────────────────────────────────────────────────────
\subsection{Detection Limitations}
% ─────────────────────────────────────────────────────────
 
\paragraph{Low recall on ID numbers (F1=0.096).}
ID number detection exhibits the largest gap between
precision and recall in our evaluation (P=0.381,
R=0.055).
The model correctly classifies an ID number when it
generates a detection, but fails to generate detections
on 94.5\% of cards.
We attribute this to two factors: (i)~ID number fields
use small fonts (8--10pt) that fall near the resolution
threshold of the EasyOCR pipeline, producing fragmented
OCR tokens that the VLM cannot confidently merge into a
single field; and (ii)~the six country-specific ID
number formats (e.g., UK National Insurance, Indian
Aadhaar, US Social Security) exhibit high visual
variance, requiring more diverse training examples than
the 2,000-card dataset provides.
Future work will address this by expanding the training
set with 500 additional cards targeting ID number field
variation, and by adding a dedicated regex sweep
post-processing step that applies country-specific
patterns as a recall safety net.
 
\paragraph{Face detection not captured in evaluation.}
Face bounding boxes are produced by a YOLOv8
object detection model operating independently of the
VLM pipeline.
Our evaluation framework measures PII detection via the
\texttt{redactions} field of the API response, which
does not include YOLOv8 face boxes.
Consequently, the evaluation reports faces F1=0.000
despite the system correctly redacting face regions
in the live application.
A future evaluation pass will query face detections
separately via a dedicated endpoint and evaluate IoU
against ground truth face bounding boxes.
 
\paragraph{Address detection precision (P=0.113).}
Address fields frequently span multiple text lines.
Our ground truth annotations record the full
multi-line bounding box, while the VLM typically
generates a single-line detection.
IoU between a one-line prediction and a three-line
ground truth box is structurally low even when
the detection is spatially correct.
Future work will adopt a containment-based matching
criterion for address fields (predicted box covers
$\geq$50\% of ground truth area) in place of strict
IoU.
 
\paragraph{Date of birth over-firing (FP rate 48\%).}
The ContextAgent correctly identifies date-of-birth
fields in most cases, but also flags expiry dates,
issue dates, and reference numbers that contain
date-like patterns.
The confusion matrix confirms 61 false positives
against 65 true positives for the DOB category.
Adding explicit exclusion rules for common non-PII
date labels (``Valid Until'', ``Issue Date'',
``Expiry'') to the ContextAgent heuristics is
expected to reduce FP rate substantially.
 
% ─────────────────────────────────────────────────────────
\subsection{Dataset and Generalization Limitations}
% ─────────────────────────────────────────────────────────
 
\paragraph{Synthetic-only training and evaluation data.}
Both the training set (2,000 cards) and the held-out
test set (150 cards) are synthetically generated using
Faker and Pillow.
Synthetic data lacks the visual noise, printing
artefacts, camera distortion, and handwritten
annotations present in real identity documents.
Transfer performance to real documents is unknown and
likely lower than the F1 scores reported here.
We identify testing on a curated set of 20--50
real anonymised identity documents as a high-priority
future step; preliminary real-document experiments will
be added before venue submission.
 
\paragraph{Six-country English-centric evaluation.}
The current system supports six countries
(Australia, Canada, Germany, India, UK, USA) with
English-language field labels.
German-language labels on German ID cards (e.g.,
\textit{Geburtsdatum} for date of birth) are
handled by the LABEL\_TO\_CATEGORY mapping but
not systematically evaluated.
Multilingual expansion using locale-specific Faker
data and multilingual OCR is planned for a future
version.
 
\paragraph{Per-country performance variance.}
Macro F1 ranges from 0.109 (Canada) to 0.214
(Germany) across countries.
Canada also accounts for three of the four pipeline
errors (500 responses), attributable to VRAM exhaustion
on complex card layouts.
These errors represent 2.0\% of the test set and
are resolved by restarting models between requests;
a future memory management improvement will implement
proactive VRAM reclamation between pipeline runs.
 
% ─────────────────────────────────────────────────────────
\subsection{Cryptographic Limitations}
% ─────────────────────────────────────────────────────────
 
\paragraph{Hash linkability without salting.}
SHA-256 is deterministic: identical documents
anchored multiple times produce the same on-chain
hash.
An adversary monitoring the blockchain can infer that
a specific document was processed repeatedly by
observing repeated hash values.
Analysis of our audit log reveals 260 duplicate hashes
among 424 records, all arising from re-processing the
same test documents during evaluation.
The mitigation is straightforward: compute the
anchor hash as SHA-256(document\_bytes $\|$ timestamp
$\|$ owner\_address) rather than SHA-256(document\_bytes)
alone.
This countermeasure has been implemented: each anchor
is salted with the owner address and a Unix timestamp,
ensuring identical documents produce distinct on-chain
hashes. The fix was validated empirically --- two runs
of the same document produced different anchor hashes
while preserving the original SHA-256 for user verification.
 
\paragraph{Ethereum mainnet gas cost.}
Each \texttt{addRecord()} call consumes 264,333 gas
units, equivalent to approximately \$13.22 USD on
Ethereum mainnet at \$2,500 per ETH.
This is economically impractical for high-volume
document processing.
Deployment on Polygon L2 reduces the per-anchor cost
to \$0.006 USD, which is viable for enterprise
compliance workflows.
The smart contract is EVM-compatible and requires
only an environment variable change to migrate from
Ganache to Polygon or Sepolia testnet.
 
\paragraph{ZK-SNARK proof generation latency.}
Groth16/BN128 proof generation via ZoKrates requires
20.3 seconds (mean) on an RTX~5060 8GB VRAM, while
on-chain verification completes in 22.1 milliseconds
(mean), confirming that verification cost is negligible.
Proof generation cost is absorbed by executing it
asynchronously in a background thread; the user
receives their redacted document before proof
generation completes.
Proof availability is therefore delayed relative to
the redaction result.
Future work will investigate STARKs or PLONK-based
schemes that offer faster proving times at the cost
of larger proof sizes, and will explore GPU-accelerated
proving libraries such as \texttt{rapidsnark}.
 
\paragraph{Trusted setup requirement.}
Groth16 proofs require a one-time trusted setup
ceremony to generate the proving and verification keys.
The current implementation uses a locally generated
setup (no multi-party ceremony), which introduces a
trust assumption: the party who ran the setup could
theoretically generate false proofs.
A production deployment should replace the local setup
with a public multi-party computation (MPC) ceremony
or migrate to a transparent proof system (STARK,
Bulletproofs) that eliminates the trusted setup
requirement entirely.
 
% ─────────────────────────────────────────────────────────
\subsection{System and Deployment Limitations}
% ─────────────────────────────────────────────────────────
 
\paragraph{Single-document processing only.}
The current system processes one document at a time.
Batch processing of document sets (e.g., a folder of
scanned records) requires sequential API calls from
the client.
A native batch endpoint with job queuing, progress
streaming, and a combined ZIP output is identified as
the next major feature addition.
 
\paragraph{Audio redaction replaces speech with silence.}
The audio pipeline silences detected PII segments
rather than replacing them with a spoken ``beep''
or synthesized neutral tone.
Silence can be as identifiable as the original
speech in some contexts (e.g., a two-second gap
where a name was spoken).
Future work will add optional bleep-tone injection
or text-to-speech replacement of redacted segments.
 
\paragraph{Local deployment only.}
Omni-Shield is designed for local-first operation
on consumer-grade hardware.
No cloud deployment, containerised orchestration,
or multi-user authentication is provided.
Deployment to a remote server requires a Cloudflare
tunnel or equivalent reverse proxy, and MetaMask wallet
authentication currently uses a single deterministic
Ganache account.
A production-grade deployment would add OAuth or
DID-based authentication and migrate blockchain
anchoring to a public testnet.
 
% ─────────────────────────────────────────────────────────
\subsection{Future Work Summary}
% ─────────────────────────────────────────────────────────
 
We identify the following concrete directions for
future work, ordered by expected impact:
 
\begin{enumerate}
 
  \item \textbf{Hash salting for linkability resistance.}
  \emph{Implemented.} Anchor hashes are now salted with
  owner address and Unix timestamp; validated empirically.
 
  \item \textbf{Expanded training with ID number focus.}
  Generate 500 additional synthetic cards with
  deliberate ID number format diversity across all
  six countries and retrain the QLoRA adapter.
 
  \item \textbf{Real document evaluation.}
  Collect 20--50 real anonymised identity documents
  and report transfer F1 alongside synthetic results.
 
  \item \textbf{Separate face evaluation pass.}
  Add a YOLOv8 face evaluation module to the
  evaluation framework and report face F1 correctly.
 
  \item \textbf{Polygon L2 deployment.}
  Replace Ganache with Sepolia testnet immediately;
  migrate to Polygon mainnet for production.
 
  \item \textbf{GRPO reinforcement learning fine-tuning.}
  Apply Group Relative Policy Optimisation with an
  IoU-based reward signal from the SFT checkpoint
  to push bounding box precision beyond the current
  QLoRA adapter.
 
  \item \textbf{Faster ZK proving.}
  Evaluate \texttt{rapidsnark} for GPU-accelerated
  Groth16 proving; alternatively migrate to PLONK
  for faster per-proof generation.
 
  \item \textbf{Multi-party trusted setup.}
  Replace the local ZK setup with a public MPC
  ceremony or a transparent proof system to eliminate
  the trusted setup assumption.
 
  \item \textbf{Batch processing endpoint.}
  Implement a native batch API with job queuing,
  progress streaming, and combined ZIP output.
 
  \item \textbf{Multilingual expansion.}
  Add locale-aware field label mappings and
  multilingual OCR support for non-English documents.
 
\end{enumerate}
 
